%! Algebra 2 — Unit 2 Cover Sheet
\documentclass[10pt]{article}
\usepackage{algebra2-article}
\usepackage{algebra2-boxes}

%=============================================================================
\begin{document}

% ── Full-bleed forest banner ────────────────────────────────────────────────────
% The header text is set INSIDE a tikz node anchored to the page so it can never
% drift above the paper edge (a negative \vspace here pushes it off the sheet).
\begin{tikzpicture}[remember picture, overlay]
  \fill[forest] (current page.north west)
    rectangle ([yshift=-1.70in]current page.north east);
  \node[anchor=north, align=center, inner sep=0pt, text width=\paperwidth]
    at ([yshift=-0.30in]current page.north) {%
      {\color{white}\bfseries\fontsize{30}{34}\selectfont Algebra\kern0.10em\ 2}\\[11pt]
      {\color{white}\bfseries\fontsize{19}{23}\selectfont Unit 2: Linear Functions}\\[11pt]%
    };
\end{tikzpicture}

% Push the body clear of the 1.70in banner (page top margin is 0.60in).
\vspace*{1.38in}

% ── Unit overview ─────────────────────────────────────────────────────────────
\begin{tcolorbox}[colback=forestbg,colframe=forest,boxrule=0.9pt,arc=2mm,
    left=4mm,right=4mm,top=3mm,bottom=3mm]
  \textbf{Unit Overview.}\quad
  Unit 2 opens the function-family study with the simplest family of all: the line.
  Students first \emph{read} a linear function's characteristics off a graph, then
  \emph{produce} the equation and the graph themselves. From there the unit follows
  absolute value --- first as an equation and inequality, then as a V-shaped function
  built by transforming a parent --- into piecewise and step functions, and closes by
  fitting a line to real bivariate data. The transformation lens introduced here
  ($f(x)+k$, $f(x+k)$, $f(kx)$, $kf(x)$) returns in every family that follows.
\end{tcolorbox}

\vspace{0.18in}

% ── Lessons in this unit ──────────────────────────────────────────────────────
\renewcommand{\arraystretch}{1.5}
\begin{skillbox}[Lessons in This Unit]{goldbox}
\begin{tabularx}{\linewidth}{@{} c >{\bfseries}p{0.30\linewidth} X @{}}
  \textbf{\#} & \textbf{Title} & \textbf{Focus} \\
  \hline
  0 & Characteristics of Linear Functions
    & Domain and range, intercepts and zeros, slope, increasing/decreasing and
      positive/negative intervals --- all read off a graph (A2.F.2a, c, f) \\
  \hline
  1 & Linear Functions: Slope, Forms \& Graphing
    & Slope as rate of change; slope-intercept, point-slope, and standard form;
      writing equations; graphing by transforming $y=x$ (A.F.1a--e) \\
  \hline
  2 & Absolute Value Equations \& Inequalities
    & $|x|$ as distance from zero; the two-case rule; solution sets written in set,
      interval, and number-line form (A2.EI.1a--e) \\
  \hline
  3 & Absolute Value Functions \& Transformations
    & The parent $y=|x|$ and vertex form $g(x)=a|x-h|+k$; vertex, axis of symmetry,
      and minimum/maximum (A2.F.1b--c; A2.F.2a, c, d, f) \\
  \hline
  4 & Piecewise \& Step Functions
    & Evaluating by piece; open/closed endpoints and discontinuities; the
      greatest-integer step function (A2.F.2a--c, f) \\
  \hline
  5 & Linear Regression
    & Scatterplots and association; the correlation coefficient $r$; the line of best
      fit; interpolation vs.\ extrapolation (A2.ST.2c--h) \\
  \hline
\end{tabularx}
\end{skillbox}

\vspace{0.14in}

% ── Standards covered ─────────────────────────────────────────────────────────
\renewcommand{\arraystretch}{1.25}
\begin{skillbox}[Virginia SOL Standards Addressed in Unit 2]{sky}
\begin{tabularx}{\linewidth}{@{} >{\bfseries}p{0.14\linewidth} X @{}}
  A.F.1   & \textit{(Algebra 1, reactivated)} Investigate, analyze, and compare linear
            functions algebraically and graphically; model linear relationships. \\
  \hline
  A2.EI.1 & Represent, solve, and interpret the solution to an absolute value equation
            or inequality in one variable. \\
  \hline
  A2.F.1  & Investigate, analyze, and compare the characteristics of a parent function
            and its transformations $f(x)+k$, $f(x+k)$, $f(kx)$, and $kf(x)$. \\
  \hline
  A2.F.2  & Investigate and analyze characteristics --- domain, range, zeros,
            intercepts, extrema, and intervals of increase, decrease, or constancy ---
            of linear, absolute value, and piecewise-defined functions, including
            graphs with discontinuities. \\
  \hline
  A2.ST.2 & Apply the data cycle with a focus on representing bivariate data with
            scatterplots and curves of best fit, and using the correlation coefficient
            to judge goodness of fit and make predictions. \\
\end{tabularx}
\end{skillbox}

\end{document}
